\documentclass[conference]{IEEEtran}

% -------------------------------------------------------
% Paquetes
% -------------------------------------------------------
\usepackage[spanish]{babel}
\usepackage[utf8]{inputenc}
\usepackage[T1]{fontenc}
\usepackage{amsmath}
\usepackage{amssymb}
\usepackage{graphicx}
\usepackage{booktabs}
\usepackage{hyperref}
\usepackage{url}
\usepackage{fancyhdr}

% -------------------------------------------------------
\begin{document}

\pagestyle{fancy}
\fancyhf{}
\rhead{\includegraphics[height=1.2cm]{assets/uts_logonegro.png}}
\fancyfoot[C]{\thepage}
\renewcommand{\headrulewidth}{0.0pt}
\setlength{\headheight}{1.5cm}

% -------------------------------------------------------
% CABECERA
% -------------------------------------------------------
\title{Acerca del Boletín | Centro de Inteligencia Artificial e Innovación\\
       \large Documento institucional del proyecto BCIAI}

\author{
  \IEEEauthorblockN{Miguel Ángel Osorio Orduz}
  \IEEEauthorblockN{Jhoan Sebastian Torres Garcia}
  \IEEEauthorblockN{Frank Camilo Zuñiga Vargas}
  \IEEEauthorblockA{Docente orientador: Carlos Pineda}
  \IEEEauthorblockA{
    Centro de Inteligencia Artificial e Innovación (CIAI)\\
    Unidades Tecnológicas de Santander --- Colombia\\
    mangelosorio@uts.edu.co\\
    jhoanstorres@uts.edu.co\\
    fczuniga@uts.edu.co
  }
}

\maketitle
\thispagestyle{fancy}

% -------------------------------------------------------
% RESUMEN
% -------------------------------------------------------
\begin{abstract}
Este documento presenta el Boletín del Centro de Inteligencia Artificial
e Innovación (BCIAI): sus objetivos, el contexto que justifica su
existencia en un momento de aceleración sin precedentes en el campo de
los modelos de lenguaje de gran escala (LLMs), el equipo docente y
estudiantil que lo sostiene, y la justificación técnica del uso de
\LaTeX{}, Bib\TeX{} y archivos \texttt{.tex} como infraestructura de
publicación. También se describe la organización del repositorio público
que aloja cada edición y las condiciones de distribución bajo la Licencia
Pública General GNU versión~3 (GPL-3.0).
\end{abstract}

\begin{IEEEkeywords}
CIAI, boletín, inteligencia artificial, LLMs, modelos de lenguaje,
\LaTeX{}, Bib\TeX{}, GPL-3.0, publicación académica, UTS
\end{IEEEkeywords}

% ===============================================================
\section{¿Qué es el Boletín CIAI?}
% ===============================================================

El Boletín del Centro de Inteligencia Artificial e Innovación (BCIAI)
es una publicación periódica producida por estudiantes del programa
de Tecnología en Desarrollo de Sistemas Informáticos de las Unidades
Tecnológicas de Santander (UTS), con el acompañamiento del cuerpo docente
del CIAI\@. Su propósito central es documentar, sintetizar y difundir
los avances más relevantes en inteligencia artificial, con énfasis en
los modelos de lenguaje de gran escala (LLMs), así como en ingeniería
de software y tecnologías afines, con una periodicidad mensual.

Cada edición cubre cuatro ejes temáticos principales: (1)~investigación
y nuevos modelos de lenguaje o sistemas de IA, incluyendo análisis
comparativos de rendimiento y costo; (2)~herramientas de código abierto
relevantes para la comunidad técnica y estudiantil; (3)~proyectos activos
dentro del CIAI que muestran el trabajo interno del centro; y
(4)~noticias del sector de ingeniería de sistemas y tecnología a nivel
global. Las ediciones se archivan de forma pública en el repositorio
oficial del proyecto~\cite{repo}, garantizando acceso libre y permanente
a cada entrega.

El boletín nació como una respuesta directa a una necesidad concreta
dentro del CIAI: contar con un canal propio de seguimiento al ecosistema
de la inteligencia artificial, que sea producido desde adentro por
quienes estudian y trabajan con estas tecnologías, y no simplemente
consumido desde fuentes externas. Esto le da un carácter formativo y
editorial que lo distingue de un simple resumen de noticias.

% ===============================================================
\section{Por qué Existe este Boletín: El Contexto de los LLMs}
% ===============================================================

La justificación más profunda del boletín radica en la naturaleza misma
del campo que cubre. Los modelos de lenguaje de gran escala han pasado,
en un período extraordinariamente corto, de ser objetos de investigación
especializada a convertirse en infraestructura tecnológica de uso masivo.
Este ritmo de cambio no tiene precedentes en la historia reciente de la
informática y plantea un desafío serio para cualquier institución
académica o profesional que quiera mantenerse al corriente.

% -------------------------------------------------------
\subsection{La velocidad de cambio en el ecosistema de LLMs}

Hasta 2022, los modelos de lenguaje eran fundamentalmente herramientas
de investigación, accesibles solo a través de interfaces técnicas y
conocidas principalmente dentro de la comunidad de aprendizaje automático.
La publicación de ChatGPT en noviembre de ese año marcó una inflexión:
en menos de dos meses alcanzó 100 millones de usuarios, convirtiéndose
en la aplicación de consumo de más rápido crecimiento en la historia.
Desde entonces, el ritmo de lanzamientos no ha hecho más que acelerarse.

En la actualidad, los laboratorios líderes ---OpenAI, Anthropic, Google,
Meta, Alibaba (Qwen), Mistral, DeepSeek, entre otros--- publican modelos
nuevos con una frecuencia de semanas o incluso días. Cada lanzamiento
trae cambios sustanciales: nuevas capacidades multimodales, ventanas de
contexto más amplias, mejoras en razonamiento, reducciones de costo por
token, arquitecturas de agentes autónomos o nuevas APIs especializadas.
Un modelo considerado de frontera en un mes puede quedar superado antes
de que termine el siguiente.

Esta velocidad de cambio tiene consecuencias directas para los
profesionales de sistemas informáticos: las herramientas, flujos de
trabajo y mejores prácticas asociadas a los LLMs cambian con una
frecuencia que ningún programa de estudios tradicional puede absorber
en tiempo real. El boletín nació precisamente para cubrir esa brecha:
ofrecer un registro mensual, curado y analizado, de lo más relevante
que ocurre en este ecosistema.

% -------------------------------------------------------
\subsection{La necesidad de un registro sistemático}

El volumen de información generada en torno a los LLMs es enorme y está
disperso. Cada laboratorio publica sus anuncios en su propio blog,
repositorio o red social; los benchmarks independientes se actualizan
continuamente en plataformas como Artificial Analysis o LMSYS Chatbot
Arena; las noticias de la industria se distribuyen entre decenas de
publicaciones especializadas. Sin un proceso de síntesis, el seguimiento
se vuelve caótico e ineficiente.

El boletín cumple la función de un registro histórico estructurado: cada
edición captura el estado del ecosistema en un momento dado, con sus
hitos, sus datos comparativos y sus eventos relevantes. Con el tiempo,
el conjunto de ediciones forma un archivo que permite trazar la evolución
del campo con precisión, sin depender de la memoria individual ni de
búsquedas ad hoc en internet. Este archivo es especialmente valioso
para trabajos de investigación, monografías o análisis de tendencias que
requieran citar hechos concretos ocurridos en fechas específicas.

% -------------------------------------------------------
\subsection{LLMs, herramientas y proyectos: tres dimensiones del mismo ecosistema}

Los modelos de lenguaje no existen en aislamiento. Alrededor de cada
avance en LLMs florecen herramientas que los aprovechan, proyectos que
los integran y debates sobre sus implicaciones. El boletín cubre estas
tres dimensiones de forma integrada porque entender solo los modelos,
sin conocer las herramientas que los rodean ni los proyectos que los
aplican, produce una visión incompleta del ecosistema.

La sección de herramientas cubre editores, IDEs, entornos de desarrollo,
aplicaciones creativas y utilidades de código abierto que emergen como
respuesta directa o indirecta al avance de la IA\@. La sección de
proyectos CIAI conecta los aprendizajes globales con el trabajo local del
centro. Y la sección de noticias contextualiza los avances técnicos con
los movimientos empresariales, regulatorios y de seguridad que definen
el entorno en que estos sistemas se despliegan.

Esta integración temática convierte al boletín en algo más que un
agregador de noticias: es un ejercicio de análisis contextual mensual
que entrena a quienes lo producen ---y a quienes lo leen--- para
entender el campo como un sistema, no como una colección de
anuncios aislados.

% -------------------------------------------------------
\subsection{El valor formativo de producir el boletín}

Desde la perspectiva pedagógica, el proceso de producción del boletín
es en sí mismo un ejercicio formativo de alta intensidad. Cada ciclo
editorial exige a los estudiantes involucrados identificar fuentes
primarias confiables, leer documentación técnica en inglés, extraer
los puntos sustanciales de un anuncio o paper, contrastar datos entre
plataformas de evaluación independientes, sintetizar hallazgos en un
formato de comunicación técnica riguroso, y manejar una cadena de
herramientas profesionales que incluye \LaTeX{}, Git y un flujo de
revisión colaborativa.

Ninguna de estas habilidades es trivial, y el hecho de practicarlas
de forma recurrente, con una fecha límite real y un producto público
como resultado, genera una curva de aprendizaje acelerada que difícilmente
se obtiene en un ambiente exclusivamente teórico. El boletín es,
en este sentido, un proyecto de formación continua disfrazado de
publicación mensual.

% ===============================================================
\section{Fin y Propósito del Boletín}
% ===============================================================

El boletín persigue tres objetivos complementarios. En primer lugar,
fortalecer las competencias investigativas de los estudiantes mediante la
búsqueda, lectura crítica y síntesis de fuentes primarias y secundarias
sobre temas de frontera tecnológica. En segundo lugar, generar un registro
histórico de los hitos de la industria de IA filtrado bajo criterios
académicos, que pueda servir como insumo bibliográfico para trabajos
futuros de la comunidad UTS\@. En tercer lugar, proyectar el trabajo
del CIAI hacia la comunidad universitaria y el público general, mostrando
los avances de los proyectos internos y posicionando al centro como un
espacio activo de innovación tecnológica dentro de Bucaramanga y la
región Santandereana.

El boletín no pretende reemplazar publicaciones científicas arbitradas;
su objetivo es divulgativo y formativo, constituyendo un ejercicio
práctico de comunicación técnica con estándares profesionales. Tampoco
pretende ser exhaustivo: el criterio editorial es la relevancia para
una audiencia de estudiantes y profesionales de sistemas informáticos,
no la cobertura total de todo lo publicado en el campo.

% ===============================================================
\section{Equipo Docente}
% ===============================================================

El proceso editorial y la orientación temática están a cargo del cuerpo
docente del CIAI en la UTS\@. El docente principal responsable de la
supervisión del boletín es:

\begin{itemize}
  \item \textbf{Carlos Pineda} --- Docente orientador del Centro de
        Inteligencia Artificial e Innovación, UTS\@. Responsable de la
        dirección académica del proyecto, la validación de contenidos y la
        aprobación de cada edición antes de su publicación.
\end{itemize}

El equipo docente define los lineamientos editoriales, garantiza la
pertinencia académica de los temas seleccionados y acompaña el proceso
formativo de los integrantes estudiantiles que producen el documento.
Su rol no es el de redactor sino el de orientador: valida que los
contenidos cumplan con estándares de rigor y que el proceso de producción
sea coherente con los objetivos pedagógicos del CIAI\@.

% ===============================================================
\section{Equipo Estudiantil}
% ===============================================================

El contenido investigativo y editorial de cada edición es producido por
los siguientes integrantes del CIAI, estudiantes del programa Tecnología
en Desarrollo de Sistemas Informáticos de la UTS:

\begin{itemize}
  \item \textbf{Miguel Ángel Osorio Orduz}
        (\texttt{mangelosorio@uts.edu.co}) --- Investigación de nuevos
        modelos de lenguaje, análisis de rendimiento comparativo y cobertura
        de lanzamientos de los principales laboratorios de IA\@.

  \item \textbf{Jhoan Sebastian Torres Garcia}
        (\texttt{jhoanstorres@uts.edu.co}) --- Cobertura de herramientas de
        código abierto, noticias del sector tecnológico y seguimiento de
        proyectos relevantes en GitHub y la comunidad de desarrollo.

  \item \textbf{Frank Camilo Zuñiga Vargas}
        (\texttt{fczuniga@uts.edu.co}) --- Proyectos internos del CIAI,
        integración y mantenimiento del repositorio del boletín, e
        infraestructura de publicación en \LaTeX{} y Git.
\end{itemize}

Los integrantes rotan responsabilidades entre ediciones con el fin de
que cada uno adquiera experiencia en todos los ejes temáticos del
boletín. Esta rotación es intencional: evita la especialización rígida
y garantiza que el equipo completo comprenda el ecosistema cubierto
desde múltiples ángulos.

% ===============================================================
\section{Justificación del Uso de \LaTeX{}, Bib\TeX{} y Archivos \texttt{.tex}}
% ===============================================================

La decisión de utilizar \LaTeX{} como sistema de composición tipográfica
responde a criterios técnicos, académicos y de mantenibilidad que se
detallan a continuación.

% -------------------------------------------------------
\subsection{\LaTeX{} como sistema de composición}

\LaTeX{} es un sistema de preparación de documentos basado en el motor
tipográfico \TeX{}, desarrollado originalmente por Donald Knuth y extendido
por Leslie Lamport~\cite{lamport}. A diferencia de los procesadores de
texto tipo \textit{What You See Is What You Get} (WYSIWYG), \LaTeX{}
separa el contenido semántico de su presentación visual, lo que garantiza
una tipografía consistente y de alta calidad sin intervención manual en
cada elemento del diseño.

Para el boletín, el uso de \LaTeX{} aporta ventajas concretas: control
preciso sobre la maquetación a doble columna (provista por la clase
\texttt{IEEEtran}), renderizado correcto de expresiones matemáticas y
tablas técnicas, compatibilidad nativa con sistemas de control de
versiones como Git (los archivos \texttt{.tex} son texto plano y
producen diffs legibles), y reproducibilidad total de la salida
tipográfica en cualquier entorno donde se disponga del compilador.

A diferencia de un archivo \texttt{.docx}, cuyo resultado visual depende
de la versión del procesador instalado, un archivo \texttt{.tex} compilado
con la misma versión del motor produce un PDF bit a bit idéntico en
cualquier sistema operativo. Esta propiedad es fundamental para un archivo
histórico que se pretende accesible y reproducible a largo plazo.

% -------------------------------------------------------
\subsection{Bib\TeX{} para la gestión de referencias}

Bib\TeX{} es el sistema estándar de gestión bibliográfica asociado a
\LaTeX{}~\cite{lamport}. Las referencias bibliográficas se almacenan en
archivos con extensión \texttt{.bib} como registros estructurados con
campos definidos (autor, título, año, URL, etc.), separados del texto
principal. El compilador resuelve automáticamente las citas dentro del
documento, genera la sección de referencias con el estilo indicado
(en las ediciones del boletín, \texttt{IEEEtran}) y garantiza que solo
aparezcan citadas las entradas que fueron referenciadas en el texto.

Esta arquitectura facilita la reutilización de referencias entre ediciones:
si un modelo o proyecto es cubierto en varias entregas, su entrada
bibliográfica se define una sola vez y se referencia desde múltiples
documentos. También facilita la corrección centralizada de datos y la
adopción de nuevos estilos de citación sin modificar el cuerpo del texto.

En ediciones donde la cantidad de referencias lo justifique, el archivo
\texttt{referencias.bib} se mantiene separado del \texttt{main.tex},
como ocurre en la edición de octubre de 2026~\cite{ediciones}. Para
documentos institucionales con pocas citas, como este, las referencias
se integran directamente con el entorno \texttt{thebibliography},
simplificando la compilación.

% -------------------------------------------------------
\subsection{Organización en archivos \texttt{.tex} modulares}

La división del documento en múltiples archivos \texttt{.tex} mediante
las directivas \verb|\input{}| o \verb|\include{}| permite separar
secciones, figuras y configuraciones en unidades independientes que se
versionan de forma granular en el repositorio~\cite{repo}. Esto es
especialmente útil cuando distintos integrantes del equipo trabajan
simultáneamente en secciones diferentes: cada uno modifica su archivo
sin generar conflictos de fusión (\textit{merge conflicts}) en el
archivo principal.

Por ejemplo, en una edición con cuatro secciones y tres autores, cada
autor puede trabajar en su propio archivo \texttt{investigacion.tex},
\texttt{herramientas.tex} o \texttt{noticias.tex}, mientras el
\texttt{main.tex} actúa únicamente como ensamblador mediante
\verb|\include{}|. El historial de Git refleja con precisión quién
modificó qué sección y cuándo, sin que los cambios en una sección
afecten el versionado de las demás.

Para ediciones sencillas o documentos institucionales como este, se
opta por un único archivo \texttt{main.tex} que concentra todo el
contenido, simplificando la compilación y la revisión.

% -------------------------------------------------------
\subsection{Ventajas frente a alternativas}

Frente a herramientas como Microsoft Word o Google Docs, \LaTeX{} ofrece
estabilidad en la maquetación cuando el documento supera unas pocas
páginas, compatibilidad nativa con clases tipográficas de conferencias y
revistas (como \texttt{IEEEtran}), y trazabilidad completa de cambios a
través de Git. El costo de entrada es mayor debido a la curva de
aprendizaje de la sintaxis, pero queda justificado en un contexto
académico donde la calidad tipográfica y la reproducibilidad son
requisitos y no opcionales.

En el contexto específico del CIAI, el uso de \LaTeX{} tiene además un
valor formativo adicional: introduce a los estudiantes en el estándar de
facto de la publicación científica en matemáticas, física, ciencias de
la computación e ingeniería. Dominar \LaTeX{} es una habilidad directamente
transferible a la producción de artículos para conferencias y revistas
arbitradas, lo que amplía las posibilidades académicas de quienes
participan en el boletín.

% ===============================================================
\section{Repositorio y Estructura de Ediciones}
% ===============================================================

El código fuente de cada edición del boletín se aloja en el repositorio
público \texttt{mikel-btw/bciai} en GitHub~\cite{repo}. La carpeta
\texttt{ediciones/} agrupa las ediciones por año y mes con la convención
\texttt{AAAA-MM/}; dentro de cada carpeta se encuentran el archivo
\texttt{main.tex} y, cuando aplica, el archivo \texttt{referencias.bib}
junto con los activos gráficos necesarios para compilar el
documento~\cite{ediciones}.

Este esquema tiene varias propiedades deseables. Primero, cada edición
es autosuficiente: contiene todo lo necesario para compilar el PDF sin
depender de archivos en otras carpetas, lo que facilita la descarga y
compilación local de ediciones específicas. Segundo, el historial de Git
conserva la autoría de cada cambio a nivel de línea, lo que hace
rastreable la contribución de cada integrante. Tercero, el repositorio
es público desde el inicio, lo que permite que cualquier persona de la
comunidad UTS o externa pueda leer, clonar, reportar errores o proponer
mejoras mediante \textit{issues} y \textit{pull requests}.

El flujo de trabajo estándar consiste en que cada integrante trabaja en
una rama (\textit{branch}) propia, genera su contenido y abre un
\textit{pull request} hacia la rama principal, donde el docente
orientador y los demás integrantes realizan la revisión antes de
fusionar. Este flujo es idéntico al utilizado en proyectos de software
profesional, lo que refuerza las competencias de trabajo colaborativo
en equipo distribuido.

% ===============================================================
\section{Licencia}
% ===============================================================

Todo el contenido de este repositorio ---textos, archivos \texttt{.tex},
archivos \texttt{.bib} y activos gráficos propios--- se distribuye bajo
la Licencia Pública General GNU versión 3.0 (GPL-3.0)~\cite{licencia}.

La GPL-3.0 garantiza las cuatro libertades fundamentales del software
libre: usar el material con cualquier propósito, estudiar cómo funciona,
modificarlo para adaptarlo a necesidades propias, y redistribuir copias
o versiones modificadas. La condición central de la licencia es el
\textit{copyleft}: cualquier trabajo derivado debe distribuirse bajo los
mismos términos, lo que asegura que las ediciones futuras y las
adaptaciones del boletín mantengan su carácter abierto y no puedan ser
apropiadas bajo términos restrictivos.

Esta elección es coherente con el espíritu del CIAI de fomentar el
conocimiento libre y accesible dentro de la comunidad universitaria y
más allá de ella. También es consistente con el enfoque editorial del
boletín, que prioriza el software de código abierto en su sección de
herramientas y valora la transparencia como principio operativo.

La licencia completa se encuentra en el archivo \texttt{LICENSE} en la
raíz del repositorio~\cite{licencia}.

% -------------------------------------------------------
\begin{thebibliography}{9}

\bibitem{repo}
  mikel-btw,
  \textit{BCIAI --- Boletín del Centro de Inteligencia Artificial e
  Innovación},
  GitHub, 2026.
  Disponible en: \url{https://github.com/mikel-btw/bciai}

\bibitem{ediciones}
  mikel-btw,
  \textit{Carpeta de ediciones del repositorio BCIAI},
  GitHub, 2026.
  Disponible en:
  \url{https://github.com/mikel-btw/bciai/tree/main/ediciones}

\bibitem{licencia}
  Free Software Foundation,
  \textit{GNU General Public License, version 3},
  FSF, 2007.
  Disponible en: \url{https://www.gnu.org/licenses/gpl-3.0.html}.
  Texto del repositorio:
  \url{https://github.com/mikel-btw/bciai/blob/main/LICENSE}

\bibitem{lamport}
  L.~Lamport,
  \textit{\LaTeX: A Document Preparation System},
  2nd~ed.
  Addison-Wesley, 1994.

\end{thebibliography}

\end{document}
